
\linksection{parametrize}{Parametrization}

\begin{minted}[autogobble]{python}
    @pytest.mark.parametrize("a, b, expected", [
        (1, 2, 3),
        (2, 3, 5),
    ])
    def test_add(a, b, expected):
        assert add(a, b) == expected
\end{minted}

Test runs twice: 1 + 2 = 3 and 2 + 3 = 5.

\linksubsection{mathspp-param}{Permutations}
\begin{minted}[autogobble]{python}
    @pytest.mark.parametrize("a", [1, 2])
    @pytest.mark.parametrize("b", [3, 4])
    def test_permutations(a, b):
        assert add(a, b) == a + b
\end{minted}

Test runs 4 times: 1 + 3 \quad 1 + 4 \quad 2 + 3 \quad 2 + 4.

Can be useful for bool flags (e.g. \py{strict=True/False} where all cases should behave the same).

\vspace{-1em}
\linksubsection{param}{pytest.param}

Instead of tuples, \py{pytest.param} objects can be passed to \mono{parametrize}, which additionally take \mono{marks} and \mono{id}:

\begin{minted}[highlightlines=2-6,escapeinside=||]{python}
@pytest.mark.parametrize("a, b, op, res", [
    pytest.param(
        2, 3, "**", 8,
        marks=pytest.mark.xfail(reason="..."),
        |id|="pow",
    ),
    pytest.param(1, 2, "+", 3, |id|="add"),
    pytest.param(3, 1, "-", 2, |id|="sub"),
])
def test_calc(a, b, op, res):
    assert calc(a, b, op) == res
\end{minted}

\begin{tabular}{ll}
    \textsb{Before:}                                             & \textsb{After:}                                                                            \\
    \mono{test_calc[2-3-**-8]} (\textcolor{termred}{\texttt{F}}) & \monot{test\_calc[\highlightbox{pow}]} \highlightbox{(\textcolor{termyellow}{\texttt{x}})} \\[.5em]
    \mono{test_calc[1-2-+-3]}                                    & \monot{test\_calc[\highlightbox{add}]}                                                     \\[.3em]
    \mono{test_calc[3-1---2]}                                    & \monot{test\_calc[\highlightbox{sub}]}
\end{tabular}

